\chapter{Differentiaalvormen en de stelling van
Stokes}\label{ch:b3:forms}

Eén stelling van de analyse heeft in twee eeuwen alle andere van
haar soort opgeslorpt: de hoofdstelling van de
integraalrekening, Green--Riemann (bewezen in het volume van
bachelorjaar 2), de divergentiestelling van Gauss, de
rotatiestelling van Kelvin--Stokes --- elk zegt dat de integraal
van een of andere afgeleide over een gebied gelijk is aan de
integraal van het oorspronkelijke object over de \omterm{def:b3:forms:boundary}{rand}. De taal
van de \emph{\omterm{def:b3:forms:diffform}{differentiaalvormen}} maakt er één uitspraak van,
$\int_M\dd\omega = \int_{\partial M}\omega$, en maakt die
uitspraak in één klap bewijsbaar. Dit hoofdstuk bouwt die taal
eerlijk op --- alternerende multilineaire algebra, de \omterm{def:b3:forms:d}{uitwendige
afgeleide}, \omterm{def:b3:forms:pullback}{terugtrekkingen}, \omterm{def:b3:forms:orientation}{oriëntatie}, integratie op de
deelvariëteiten van \cref{ch:b3:submanifolds} --- bewijst de
stelling van Stokes, en int de eerste cheques: de klassieke
integraalstellingen, het \omterm{def:b3:forms:winding}{windingsgetal} dat in het geheim
\cref{ch:b3:residues} aandreef, en, in de weekendopgave, de
vastepuntstelling van Brouwer. Overal betekent \emph{glad}
$\mathcal C^\infty$; elke afbeelding en elke vorm is glad tenzij
anders vermeld. Dat kost geen algemeenheid die het waard is, en
maakt de handen vrij.

\section{Alternerende multilineaire algebra}

\begin{definition}\label{def:b3:forms:alternating}
Zij $E$ een reële vectorruimte van dimensie $n$. Een
\emph{$k$-lineaire alternerende vorm}\index{alternerende vorm}
op $E$ is een afbeelding $\alpha\colon E^k \to \R$, lineair in
elke veranderlijke, met $\alpha(v_1, \dots, v_k) = 0$ zodra twee
argumenten gelijk zijn. Hun ruimte noteren we $\Lambda^k E^*$;
per afspraak is $\Lambda^0E^* = \R$. De alternantie dwingt de
antisymmetrie af: twee argumenten verwisselen verandert het teken
(ontwikkel $\alpha(\dots, v + w, \dots, v + w, \dots) = 0$), en
algemener is $\alpha(v_{\sigma(1)}, \dots, v_{\sigma(k)}) =
\varepsilon(\sigma)\,\alpha(v_1, \dots, v_k)$ voor elke
permutatie $\sigma$.
\end{definition}

\begin{example}
Op $E = \R^n$ is $\Lambda^1E^* = E^*$ de \omterm{def:b3:banach:operator}{duale ruimte}; de
determinant in de canonieke basis is een $n$-lineaire
\omterm{def:b3:forms:alternating}{alternerende vorm}, en \cref{prop:b3:forms:basis} zal tonen dat ze
$\Lambda^nE^*$ opspant --- de diepe reden waarom de determinant
op een schaalfactor na uniek is. Voor $k > n$ is $\Lambda^kE^* =
\{0\}$: $k$ vectoren zijn afhankelijk, en er een langs de andere
ontwikkelen doodt $\alpha$ wegens de alternantie.
\end{example}

\begin{definition}\label{def:b3:forms:wedge}
Voor $\ell_1, \dots, \ell_k \in E^*$ is hun \emph{uitwendig
product}\index{uitwendig product} de $k$-lineaire \omterm{def:b3:forms:alternating}{alternerende
vorm}
\[
(\ell_1 \wedge \dots \wedge \ell_k)(v_1, \dots, v_k)
= \det\bigl(\ell_i(v_j)\bigr)_{1 \leq i, j \leq k} .
\]
De alternantie en de multilineariteit zijn die van de determinant
in haar kolommen.
\end{definition}

\begin{proposition}[Basis van $\Lambda^k$]\label{prop:b3:forms:basis}
Zij $(e_1, \dots, e_n)$ een basis van $E$ met duale basis
$(e_1^*, \dots, e_n^*)$. De vormen
\[
e_I^* = e_{i_1}^* \wedge \dots \wedge e_{i_k}^*,
\qquad I = \{i_1 < \dots < i_k\} \subseteq \{1, \dots, n\},
\]
vormen een basis van $\Lambda^kE^*$; dus is $\dim\Lambda^kE^* =
\binom nk$. Expliciet is $\alpha = \sum_{\abs I =
k}\alpha(e_{i_1}, \dots, e_{i_k})\,e_I^*$.
\end{proposition}

\begin{proof}
\emph{Voortbrengend.} Zij $\alpha \in \Lambda^kE^*$ en $\beta =
\sum_I\alpha(e_I)\,e_I^*$, waarbij $\alpha(e_I)$ afkort voor
$\alpha(e_{i_1}, \dots, e_{i_k})$. Beide leden zijn $k$-lineair
en alternerend, dus vallen ze samen zodra ze samenvallen op alle
$k$-tallen $(e_{j_1}, \dots, e_{j_k})$ met $j_1 < \dots < j_k$
(de multilineariteit herleidt tot $k$-tallen basisvectoren, de
alternantie tot strikt stijgende). En $e_I^*(e_{j_1}, \dots,
e_{j_k}) = \det(e_{i_r}^*(e_{j_s})) = \delta_{IJ}$: voor $I = J$
is de matrix de identiteit; voor $I \neq J$ is een rij nul. Dus
$\beta(e_J) = \alpha(e_J)$ voor alle $J$: $\beta = \alpha$.
\emph{Vrijheid.} Is $\sum_I c_Ie_I^* = 0$, dan geeft evalueren op
$(e_{j_1}, \dots, e_{j_k})$ dat $c_J = 0$.
\end{proof}

\begin{definition}\label{def:b3:forms:wedgegeneral}
Het \omterm{def:b3:forms:wedge}{uitwendig product} breidt uit tot een bilineaire afbeelding
$\Lambda^kE^* \times \Lambda^\ell E^* \to \Lambda^{k+\ell}E^*$,
bepaald door de bilineariteit en $(e_I^*) \wedge (e_J^*) = e_I^*
\wedge e_J^*$ (aaneenschakelen en herordenen; het product is $0$
als $I \cap J \neq \varnothing$). Ze is associatief en
\emph{gegradeerd anticommutatief}:
\[
\beta \wedge \alpha = (-1)^{k\ell}\,\alpha \wedge \beta
\qquad (\alpha \in \Lambda^k, \ \beta \in \Lambda^\ell).
\]
\end{definition}

\begin{proof}[\omnameProof]
Beide eigenschappen worden op basiselementen nagegaan en met de
bilineariteit uitgebreid. Associativiteit: beide manieren om
$e_I^* \wedge e_J^* \wedge e_K^*$ te haken zijn gelijk aan het
\omterm{def:b3:forms:wedge}{uitwendig product} van de aaneengeschakelde familie $1$-vormen,
volgens de determinantformule van \cref{def:b3:forms:wedge}
(ontwikkeling van Laplace per blok). De tekenregel: elk van de
$\ell$ factoren van $\beta$ langs de $k$ factoren van $\alpha$
verplaatsen kost een teken per naburige verwisseling (het
omwisselen van twee rijen van de determinant), dus in totaal
$(-1)^{k\ell}$.
\end{proof}

\begin{proposition}[Terugtrekking, lineair geval]
\label{prop:b3:forms:linearpullback}
Een lineaire afbeelding $u\colon E \to F$ induceert voor elke $k$
de lineaire afbeelding $u^*\colon \Lambda^kF^* \to \Lambda^kE^*$,
$(u^*\alpha)(v_1, \dots, v_k) = \alpha(u(v_1), \dots, u(v_k))$.
Ze voldoet aan $u^*(\alpha \wedge \beta) = u^*\alpha \wedge
u^*\beta$ en $(u \circ w)^* = w^* \circ u^*$. Bovendien:
\begin{enumerate}
  \item is $\dim E = n$ en $u\colon E \to E$, dan is op de rechte
        $\Lambda^nE^*$ $u^*\alpha = (\det u)\,\alpha$;
  \item is $\operatorname{rk}u < k$, dan is $u^* = 0$ op
        $\Lambda^kF^*$.
\end{enumerate}
\end{proposition}

\begin{proof}
De functoriële identiteiten volgen onmiddellijk uit de definities
(controleer voor de productregel op \omterm{def:b3:forms:wedge}{uitwendige producten} van
$1$-vormen met de determinantformule --- $\det(\ell_i(uv_j)) =
\det((u^*\ell_i)(v_j))$ --- en breid bilineair uit).
(1) $u^*$ beeldt de eendimensionale $\Lambda^nE^*$
(\cref{prop:b3:forms:basis}) in zichzelf af, dus is $u^*\alpha =
c\,\alpha$ met $c$ onafhankelijk van $\alpha \neq 0$; testen op
$\alpha = e_1^* \wedge \dots \wedge e_n^*$ en $(v_j) = (e_j)$
geeft $c = \det(e_i^*(ue_j)) = \det u$.
(2) Voor $v_1, \dots, v_k \in E$ liggen de vectoren $u(v_1),
\dots, u(v_k)$ in het beeld van $u$, van dimensie $< k$: ze zijn
lineair afhankelijk, en een \omterm{def:b3:forms:alternating}{alternerende vorm} verdwijnt op een
afhankelijke familie (ontwikkel de afhankelijke vector langs de
andere).
\end{proof}

\section{Differentiaalvormen en de uitwendige afgeleide}

\begin{definition}\label{def:b3:forms:diffform}
Zij $U \subseteq \R^n$ open. Een
\emph{differentiaal-$k$-vorm}\index{differentiaalvorm} op $U$ is
een gladde afbeelding $\omega\colon U \to \Lambda^k(\R^n)^*$; in
de basis van \cref{prop:b3:forms:basis} (met $\dd x_i$ voor
$e_i^*$):
\[
\omega = \sum_{\abs I = k} a_I\,\dd x_I,
\qquad
\dd x_I = \dd x_{i_1} \wedge \dots \wedge \dd x_{i_k},
\]
met gladde coëfficiënten $a_I \in \mathcal C^\infty(U)$. Hun
ruimte is $\Omega^k(U)$; $\Omega^0(U) = \mathcal C^\infty(U)$.
Een $0$-vorm is een functie; een $1$-vorm is een veld van
lineaire vormen (bijvoorbeeld de differentiaal $\dd f$ van een
functie); een $n$-vorm is $a\,\dd x_1\wedge\dots\wedge\dd x_n$,
de natuurlijke integrand van \cref{ch:b3:product}.
\end{definition}

\begin{definition}[Uitwendige afgeleide]\label{def:b3:forms:d}
De \emph{uitwendige afgeleide}\index{uitwendige afgeleide} is de
lineaire afbeelding $\dd\colon \Omega^k(U) \to \Omega^{k+1}(U)$
gedefinieerd door
\[
\dd\Bigl(\sum_I a_I\,\dd x_I\Bigr)
= \sum_I \dd a_I \wedge \dd x_I
= \sum_I\sum_{j=1}^n \frac{\partial a_I}{\partial
x_j}\,\dd x_j \wedge \dd x_I .
\]
Op $0$-vormen is ze de gewone differentiaal.
\end{definition}

\begin{theorem}\label{thm:b3:forms:dsquare}
(a) $\dd(\omega \wedge \eta) = \dd\omega \wedge \eta +
(-1)^k\,\omega \wedge \dd\eta$ voor $\omega \in \Omega^k$ (de
\emph{gegradeerde regel van Leibniz}).
(b) $\dd \circ \dd = 0$.\index{d kwadraat nul@$\dd^2 = 0$}
\end{theorem}

\begin{proof}
(a) Wegens de bilineariteit volstaat het $\omega = a\,\dd x_I$ en
$\eta = b\,\dd x_J$ te behandelen. Dan is $\omega \wedge \eta =
ab\,\dd x_I \wedge \dd x_J$ en
\[
\dd(\omega\wedge\eta) = (b\,\dd a + a\,\dd b) \wedge \dd x_I
\wedge \dd x_J
= (\dd a \wedge \dd x_I) \wedge (b\,\dd x_J) +
a\,\dd b \wedge \dd x_I \wedge \dd x_J,
\]
en de $1$-vorm $\dd b$ langs de $k$ factoren van $\dd x_I$
verplaatsen kost $(-1)^k$
(\cref{def:b3:forms:wedgegeneral}): de tweede term is
$(-1)^k\,\omega \wedge \dd\eta$.
(b) Voor een functie: $\dd(\dd a) = \sum_{i,j}
\frac{\partial^2a}{\partial x_j\partial x_i}\,\dd x_j \wedge \dd
x_i$. De coëfficiënt is symmetrisch in $(i,j)$ volgens de
stelling van Schwarz over gemengde partiële afgeleiden (bewezen
in het volume van bachelorjaar 2; $a$ is $\mathcal C^\infty$),
terwijl $\dd x_j \wedge \dd x_i$ antisymmetrisch is: de termen
$(i,j)$ en $(j,i)$ paren en alles valt weg. Voor een algemene
$\omega = \sum a_I\dd x_I$: $\dd\dd\omega = \sum\dd(\dd a_I
\wedge \dd x_I) = \sum(\dd\dd a_I\wedge\dd x_I - \dd
a_I\wedge\dd(\dd x_I))$ volgens (a), en beide termen verdwijnen
($\dd(\dd x_I) = 0$ omdat de coëfficiënt constant is).
\end{proof}

\begin{definition}[Terugtrekking]\label{def:b3:forms:pullback}
Zij $\varphi\colon U \to V$ glad ($U \subseteq \R^m$, $V
\subseteq \R^n$ open). De \emph{terugtrekking}\index{terugtrekking}
$\varphi^*\colon \Omega^k(V) \to \Omega^k(U)$ is puntsgewijs
gedefinieerd met de lineaire terugtrekking langs de
differentiaal: $(\varphi^*\omega)_x =
(D\varphi(x))^*\,\omega_{\varphi(x)}$. Concreet substitueert
$\varphi^*$: $\varphi^*f = f \circ \varphi$ op functies,
$\varphi^*(\dd y_i) = \dd\varphi_i =
\sum_j\frac{\partial\varphi_i}{\partial x_j}\dd x_j$, en
$\varphi^*(a\,\dd y_{i_1}\wedge\dots\wedge \dd y_{i_k}) =
(a\circ\varphi)\,\dd\varphi_{i_1}\wedge\dots\wedge
\dd\varphi_{i_k}$.
\end{definition}

\begin{theorem}\label{thm:b3:forms:pullback}
(a) $\varphi^*(\omega \wedge \eta) = \varphi^*\omega \wedge
\varphi^*\eta$ en $(\psi \circ \varphi)^* = \varphi^* \circ
\psi^*$.
(b) $\varphi^*(\dd\omega) = \dd(\varphi^*\omega)$: de \omterm{def:b3:forms:d}{uitwendige
afgeleide} commuteert met elke gladde substitutie --- de
identiteit die haar tot \emph{de} afgeleide van de theorie
maakt.
(c) Is $\varphi\colon U \to V$ glad tussen open delen van $\R^n$
en $\omega = a\,\dd y_1 \wedge \dots \wedge \dd y_n$, dan is
$\varphi^*\omega = (a \circ \varphi)\,\det\bigl(D\varphi\bigr)\,
\dd x_1 \wedge \dots \wedge \dd x_n$.
\end{theorem}

\begin{proof}
(a) Puntsgewijze uitspraken over lineaire \omterm{def:b3:forms:pullback}{terugtrekkingen}
(\cref{prop:b3:forms:linearpullback}), plus de kettingregel
$D(\psi\circ\varphi)(x) = D\psi(\varphi(x))\,D\varphi(x)$.
(b) Voor een $0$-vorm $f$: $\varphi^*(\dd f) = \dd f \circ
D\varphi = \dd(f \circ \varphi)$ is de kettingregel. Voor $\omega
= a\,\dd y_I$: met (a) is $\varphi^*\omega =
(a\circ\varphi)\,\dd\varphi_{i_1}\wedge\dots\wedge
\dd\varphi_{i_k}$, dus geven de regel van Leibniz
(\cref{thm:b3:forms:dsquare}(a)) en $\dd\dd\varphi_{i_r} = 0$
\[
\dd(\varphi^*\omega) = \dd(a\circ\varphi) \wedge
\dd\varphi_{i_1}\wedge\dots\wedge\dd\varphi_{i_k}
= \varphi^*(\dd a) \wedge \varphi^*(\dd y_I)
= \varphi^*(\dd a \wedge \dd y_I) =
\varphi^*(\dd\omega).
\]
(c) Puntsgewijs is dat precies $u^*\alpha = (\det u)\alpha$ op
vormen van topgraad (\cref{prop:b3:forms:linearpullback}(1) met
$u = D\varphi(x)$).
\end{proof}

\begin{example}[Poolcoördinaten]\label{ex:b3:forms:polar}
Voor $\varphi(r, \theta) = (r\cos\theta, r\sin\theta)$ is
$\varphi^*\dd x = \cos\theta\,\dd r - r\sin\theta\,\dd\theta$ en
$\varphi^*\dd y = \sin\theta\,\dd r + r\cos\theta\,\dd\theta$,
dus
\[
\varphi^*(\dd x \wedge \dd y) =
(\cos\theta\,\dd r - r\sin\theta\,\dd\theta) \wedge
(\sin\theta\,\dd r + r\cos\theta\,\dd\theta)
= r\,\dd r \wedge \dd\theta,
\]
waarbij de jacobiaan van \cref{ex:b3:product:polar} met pure
algebra opduikt --- geen maattheorie. \cref{exo:b3:forms:9}
maakt van die opmerking een uitspraak: voor georiënteerde
integralen \emph{is} de formule van de variabelensubstitutie de
terugtrekkingsformule.
\end{example}

\section{Gesloten en exacte vormen; het lemma van Poincaré}

\begin{definition}\label{def:b3:forms:closedexact}
$\omega \in \Omega^k(U)$ heet \emph{gesloten}\index{gesloten
vorm} als $\dd\omega = 0$, en \emph{exact}\index{exacte vorm} als
$\omega = \dd\eta$ voor een $\eta \in \Omega^{k-1}(U)$ (een
\emph{primitieve} van $\omega$). Exact $\Rightarrow$ gesloten
volgens $\dd^2 = 0$; de omkering is een vraag over de
\emph{vorm} van $U$.
\end{definition}

\begin{example}[De hoekvorm]\label{ex:b3:forms:angular}
Op $U = \R^2 \setminus \{0\}$ is
\[
\omega_\theta = \frac{x\,\dd y - y\,\dd x}{x^2 + y^2}
\]
\omterm{def:b3:forms:closedexact}{gesloten} (rechtstreekse berekening: \cref{exo:b3:forms:4}) maar
niet \omterm{def:b3:forms:closedexact}{exact}: haar integraal langs de eenheidscirkel is $2\pi \neq
0$, terwijl de integralen van \omterm{def:b3:forms:closedexact}{exacte vormen} langs \omterm{def:b3:forms:closedexact}{gesloten}
krommen verdwijnen (\cref{prop:b3:forms:exactloop}). Lokaal is
$\omega_\theta = \dd\theta$ voor elke gladde bepaling $\theta$
van de poolhoek --- vandaar de naam en de obstructie: zo'n
bepaling bestaat niet op heel $U$. Deze ene vorm drijft het
\omterm{def:b3:forms:winding}{windingsgetal} aan (\cref{sec:b3:forms:winding}) en, daardoorheen,
de residustelling van \cref{ch:b3:residues}.
\end{example}

\begin{theorem}[Lemma van Poincaré]\label{thm:b3:forms:poincare}
Zij $U \subseteq \R^n$ open en \emph{stervormig} ten opzichte van
$0$. Elke \omterm{def:b3:forms:closedexact}{gesloten} $k$-vorm op $U$ ($k \geq 1$) is
\omterm{def:b3:forms:closedexact}{exact}.\index{lemma van Poincaré}
\end{theorem}

\begin{proof}
We bouwen een lineaire \emph{homotopie-operator} $h\colon
\Omega^k(U) \to \Omega^{k-1}(U)$ met
\begin{equation}\label{eq:b3:forms:homotopy}
\dd(h\omega) + h(\dd\omega) = \omega
\qquad (k \geq 1);
\end{equation}
is $\dd\omega = 0$, dan is $\omega = \dd(h\omega)$ en zijn we
klaar. Stel voor $\omega = \sum_I a_I\,\dd x_I$
\[
h\omega = \sum_{\abs I = k}\ \sum_{r=1}^{k}(-1)^{r-1}
\Bigl(\int_0^1 t^{k-1}a_I(tx)\,\dd t\Bigr)\,
x_{i_r}\,\dd x_{i_1}\wedge\dots\wedge
\widehat{\dd x_{i_r}}\wedge\dots\wedge\dd x_{i_k}
\]
(het dakje schrapt een factor; de integralen zijn glad in $x$
wegens het differentiëren onder de integraal,
\cref{thm:b3:lebesgue:paramdiff}, waarbij alle afgeleiden op
\omterm{def:b3:topology:compact}{compacte} verzamelingen gedomineerd worden).
\eqref{eq:b3:forms:homotopy} nagaan is een berekening die je één
keer in je leven doet, dus we doen ze voluit. Leg $I$ vast en
neem $\omega = a\,\dd x_I$ (lineariteit). Eerst is
\[
\dd(h\omega) = k\Bigl(\int_0^1t^{k-1}a(tx)\dd t\Bigr)\dd x_I
+ \sum_{r=1}^k(-1)^{r-1}\sum_{j=1}^n
\Bigl(\int_0^1t^{k}\,\partial_ja(tx)\,\dd t\Bigr)
x_{i_r}\,\dd x_j\wedge\dd x_{I\setminus i_r} :
\]
de eerste groep verzamelt de termen waarin $\dd$ de factor
$x_{i_r}$ raakt --- het \omterm{def:b3:forms:wedge}{uitwendig product} $\dd x_{i_r}\wedge\dd
x_{I\setminus i_r}$ stelt $\dd x_I$ opnieuw samen met een teken
$(-1)^{r-1}$ dat de voorfactor opheft, en de $k$ waarden van $r$
geven de factor $k$ --- terwijl de tweede groep de termen
verzamelt waar $\dd$ de integraal raakt (de kettingregel brengt
$t\,\partial_ja(tx)$ naar buiten). Vervolgens is $\dd\omega =
\sum_j\partial_ja\,\dd x_j\wedge\dd x_I$, en de definitie van $h$
in graad $k+1$ toepassen, waarbij de index $j$ de eerste plaats
inneemt, geeft
\[
h(\dd\omega) = \sum_{j=1}^n\Bigl(\int_0^1t^{k}
\partial_ja(tx)\dd t\Bigr)x_j\,\dd x_I
- \sum_{j=1}^n\sum_{r=1}^k(-1)^{r-1}
\Bigl(\int_0^1t^{k}\partial_ja(tx)\dd t\Bigr)
x_{i_r}\,\dd x_j\wedge\dd x_{I\setminus i_r} .
\]
De dubbele sommen heffen elkaar op in $\dd(h\omega) +
h(\dd\omega)$, dat dus gelijk is aan
\[
\Bigl(\int_0^1\bigl(k\,t^{k-1}a(tx) +
t^k{\textstyle\sum_j}x_j\,\partial_ja(tx)\bigr)\dd t\Bigr)
\dd x_I
= \Bigl(\int_0^1\frac{\dd}{\dd t}\bigl(t^ka(tx)\bigr)\dd
t\Bigr)\dd x_I
= a(x)\,\dd x_I = \omega,
\]
volgens de hoofdstelling van de integraalrekening. De
stervormigheid kwam binnen waar ze moest: $tx \in U$ voor $t \in
\intcc01$, zodat $a(tx)$ zin heeft.
\end{proof}

\begin{remark}
Voor $k = 1$ en $\omega = \sum_j a_j\dd x_j$ is de primitieve
$f(x) = \int_0^1\sum_ja_j(tx)\,x_j\,\dd t$ --- de lijnintegraal
van $\omega$ langs het lijnstuk $[0, x]$: de stelling is het
``een veld met symmetrische jacobiaan is een gradiënt'' in
meerdere veranderlijken uit het volume van bachelorjaar 2, nu in
elke graad. De hoekvorm (\cref{ex:b3:forms:angular}) toont dat de
hypothese over $U$ geen versiering is: $\R^2\setminus\{0\}$ is
niet stervormig, en daar impliceert geslotenheid geen exactheid.
Wat op een algemene \omterm{def:b3:topology:topology}{open verzameling} overleeft, wordt gemeten
door de \emph{de-rhamcohomologie} $H^k(U) =
\ker\dd/\operatorname{im}\dd$ --- zie \cref{exo:b3:forms:12} voor
de eerste niet-triviale berekening.
\end{remark}

\section{Oriëntatie en integratie op deelvariëteiten}

Een $k$-vorm integreren vereist $k$-dimensionaal georiënteerd
terrein. Herinner je uit \cref{ch:b3:submanifolds}
(\cref{thm:b3:submanifolds:characterizations}) dat een
$k$-deelvariëteit $M \subseteq \R^n$ lokaal het beeld is van een
reguliere parametrisering $\gamma\colon V \to M \cap W$ ($V
\subseteq \R^k$ open, $\gamma$ een \omterm{def:b3:topology:continuity}{homeomorfisme} op haar beeld
met injectieve differentiaal).

\begin{definition}\label{def:b3:forms:orientation}
Een \emph{oriëntatie}\index{oriëntatie} van $M$ is een keuze,
voor elke $p \in M$, van een van de twee oriëntatieklassen van
bases van de \omterm{def:b3:submanifolds:tangent}{raakruimte} $T_pM$, die \emph{lokaal coherent} is:
rond elk punt is er een parametrisering $\gamma$ waarvan het
coördinaatframe $(\partial_1\gamma, \dots, \partial_k\gamma)$ in
elk punt van haar domein positief georiënteerd is. Zulke
parametriseringen heten \emph{direct}. $M$ heet
\emph{oriënteerbaar} als er een oriëntatie bestaat; de
möbiusband toont dat dat kan falen. Alle deelvariëteiten in dit
hoofdstuk zijn georiënteerd.
\end{definition}

\begin{definition}[Integraal van een vorm]\label{def:b3:forms:integral}
Zij $M$ een georiënteerde $k$-deelvariëteit en $\omega$ een
$k$-vorm gedefinieerd op een \omterm{def:b3:topology:topology}{omgeving} van $M$, met
$\operatorname{supp}\omega \cap M$ \omterm{def:b3:topology:compact}{compact}.
(a) Is $\operatorname{supp}\omega \cap M \subseteq \gamma(V)$
voor één enkele directe parametrisering, stel dan
\[
\int_M\omega = \int_V\gamma^*\omega
\]
--- waarbij het rechterlid de lebesgue-integraal over $V$ is
(\cref{ch:b3:product}) van de coëfficiënt $g$ van
$\gamma^*\omega = g\,\dd u_1\wedge\dots\wedge\dd u_k$, die
\omterm{def:b3:topology:continuity}{continu} is met \omterm{def:b3:topology:compact}{compacte} drager.
(b) Kies in het algemeen eindig veel directe parametriseringen
$\gamma_i(V_i)$ die de \omterm{def:b3:topology:compact}{compacte} $\operatorname{supp}\omega \cap
M$ overdekken en een ondergeschikte \omterm{lem:b3:forms:partition}{partitie van de eenheid}
$(\chi_i)$ (\cref{lem:b3:forms:partition}), en stel
$\int_M\omega = \sum_i\int_M\chi_i\,\omega$, waarbij elke term
met (a) wordt berekend.
Voor een kromme ($k = 1$) geparametriseerd door
$\gamma\colon\intcc ab\to\R^n$ schrijven we $\int_\gamma\omega =
\int_a^b\gamma^*\omega$, zonder dat injectiviteit vereist is.
\end{definition}

\begin{lemma}[Consistentie]\label{lem:b3:forms:consistency}
Definitie (a) hangt niet van de directe parametrisering af, en
definitie (b) hangt noch van de overdekking noch van de \omterm{lem:b3:forms:partition}{partitie
van de eenheid} af. Bovendien geldt: is $\Phi$ een diffeomorfisme
van \omterm{def:b3:topology:topology}{omgevingen} van twee georiënteerde deelvariëteiten met
$\Phi(M) = M'$, dat in een punt van elke component van $M$ een
direct frame van $M$ naar een direct frame van $M'$ brengt, dan
is $\int_{M'}\omega = \int_M\Phi^*\omega$.
\end{lemma}

\begin{proof}
(a) Zij $\gamma\colon V \to M$ en $\delta\colon V' \to M$ directe
parametriseringen waarvan de beelden $\operatorname{supp}\omega
\cap M$ bevatten. De overgang $\tau = \delta^{-1}\circ\gamma$ is
een diffeomorfisme tussen de betrokken open delen van $V$ en $V'$
(de gladheid van de overgangen:
\cref{thm:b3:submanifolds:characterizations}, via de lokale
beschrijving met grafieken), en daar is $\gamma =
\delta\circ\tau$, dus $\gamma^*\omega = \tau^*(\delta^*\omega)$
(\cref{thm:b3:forms:pullback}(a)). Schrijf $\delta^*\omega =
g\,\dd u_1\wedge\dots\wedge\dd u_k$; dan is
(\cref{thm:b3:forms:pullback}(c)) $\tau^*(\delta^*\omega) =
(g\circ\tau)\,\det(D\tau)\,\dd u_1\wedge\dots\wedge\dd u_k$.
Omdat beide frames direct zijn, beeldt $D\tau$ een positieve
basis op een positieve basis af: $\det D\tau > 0$, dus $\det
D\tau = \abs{\det D\tau}$, en de stelling over de
variabelensubstitutie (\cref{thm:b3:product:changeofvar}) geeft
$\int(g\circ\tau)\abs{\det D\tau} = \int g$: de twee integralen
vallen samen. \emph{Hier ligt de hele reden waarom \omterm{def:b3:forms:orientation}{oriëntatie}
bestaat}: zonder controle over het teken verschillen de jacobiaan
en haar absolute waarde, en is de integraal slecht gedefinieerd.
(b) Zijn $(\chi_i)$ en $(\tilde\chi_j)$ twee toelaatbare
partities (overdekkingen inbegrepen), dan is volgens (a) en de
eindige additiviteit $\sum_i\int\chi_i\omega =
\sum_{i,j}\int\chi_i\tilde\chi_j\omega =
\sum_j\int\tilde\chi_j\omega$, waarbij elke dubbele term in beide
kaarten berekend kan worden. De laatste uitspraak: doorloopt
$\gamma$ de directe parametriseringen van $M$, dan doorloopt
$\Phi\circ\gamma$ de directe parametriseringen van $M'$ (de
vergelijking van de \omterm{def:b3:forms:orientation}{oriëntaties} is lokaal constant en in één punt
per component vastgelegd), en is $(\Phi\circ\gamma)^*\omega =
\gamma^*(\Phi^*\omega)$.
\end{proof}

\begin{lemma}[Partities van de eenheid, compact geval]
\label{lem:b3:forms:partition}
Zij $K \subseteq \R^n$ \omterm{def:b3:topology:compact}{compact} en $W_1, \dots, W_m$ \omterm{def:b3:topology:topology}{open
verzamelingen} die $K$ overdekken. Er bestaan $\chi_1, \dots,
\chi_m \in \mathcal C^\infty(\R^n)$ met $0 \leq \chi_i \leq 1$,
$\operatorname{supp}\chi_i \subseteq W_i$ \omterm{def:b3:topology:compact}{compact}, en $\sum\chi_i
= 1$ op een \omterm{def:b3:topology:topology}{omgeving} van $K$.\index{partitie van de eenheid}
\end{lemma}

\begin{proof}
Elke $x \in K$ ligt in een $W_{i(x)}$ met een \omterm{def:b3:forms:closedexact}{gesloten} bal $\bar
B(x, 2r_x) \subseteq W_{i(x)}$; de \omterm{def:b3:topology:compact}{compactheid} haalt er $x_1,
\dots, x_N$ uit waarvoor de ballen $B(x_s, r_{x_s})$ $K$
overdekken. Neem voor elke $s$ een bult $\theta_s \in \mathcal
C^\infty$ met $0 \leq \theta_s \leq 1$, $\theta_s = 1$ op $\bar
B(x_s, r_{x_s})$ en $\operatorname{supp}\theta_s \subseteq B(x_s,
2r_{x_s})$ (strijk de indicator van de bal met straal
$\frac32r_{x_s}$ glad, \cref{thm:b3:lp:regularization}). Wijs
elke $s$ toe aan één index $i(s)$ met $B(x_s, 2r_{x_s})
\subseteq W_{i(s)}$ en stel $\Theta_i = \sum_{i(s) = i}\theta_s$.
Op de \omterm{def:b3:topology:topology}{open verzameling} $\Omega_0 = \{\sum_j\Theta_j >
\tfrac12\} \supseteq K$ doen de functies
$\Theta_i/\sum_j\Theta_j$ hun werk, maar ze zijn enkel daar
gedefinieerd; om te globaliseren, zij $\rho \in \mathcal
C^\infty(\R^n)$ met $\rho = 0$ waar $\sum_j\Theta_j \geq 1$ en
$\rho > 0$ waar $\sum_j\Theta_j \leq \tfrac12$ (strijk een
geschikte afknotting van $1 - \sum_j\Theta_j$ glad), en stel
\[
\chi_i = \frac{\Theta_i}{\rho + \sum_j\Theta_j} .
\]
De noemer is overal $> 0$ en gelijk aan $\sum_j\Theta_j$ op
$\{\sum_j\Theta_j \geq 1\}$, een open \omterm{def:b3:topology:topology}{omgeving} van $K$ (elk punt
van $K$ ligt in een bal $B(x_s, r_{x_s})$ waar $\theta_s = 1$);
daar is $\sum_i\chi_i = 1$. De dragers en de grenzen zijn
duidelijk.
\end{proof}

\section{De stelling van Stokes}

\begin{definition}\label{def:b3:forms:boundary}
Een \emph{$k$-deelvariëteit met rand}\index{deelvariëteit met
rand} $M \subseteq \R^n$ is een verzameling die door reguliere
parametriseringen van twee soorten wordt overdekt: \emph{inwendige
kaarten} $\gamma\colon V \to M \cap W$ met $V \subseteq \R^k$
open, en \emph{randkaarten} $\gamma\colon V \cap H^k \to M \cap
W$, waarbij $H^k = \{u \in \R^k : u_k \geq 0\}$ en $\gamma$ zich
glad en regulier uitbreidt tot de open $V$. De \emph{rand}
$\partial M$ is de verzameling punten die bij $u_k = 0$ worden
bereikt; ze is een $(k-1)$-deelvariëteit zonder rand,
geparametriseerd door de afbeeldingen $u' \mapsto \gamma(u', 0)$.
Een \omterm{def:b3:forms:orientation}{oriëntatie} van $M$ induceert er een op $\partial M$ met de
regel \emph{buitenwaartse normaal eerst}: in $p \in \partial M$
is een basis $(w_1, \dots, w_{k-1})$ van $T_p\partial M$ positief
dan en slechts dan als $(\nu, w_1, \dots, w_{k-1})$ een positieve
basis van $T_pM$ is, waarbij $\nu \in T_pM \setminus T_p\partial
M$ uit $M$ naar buiten wijst (in een randkaart: $\nu =
-\partial_k\gamma$, op het optellen van tangentiële componenten
na --- de oriëntatieklasse ziet die niet).
\end{definition}

\begin{omfigure}
\begin{tikzpicture}[scale=1.05]
  \fill[omDef!10] (-2,0) arc (180:0:2) -- cycle;
  \draw[thick, omDef] (-2,0) arc (180:0:2);
  \draw[thick, omThm] (-2,0) -- (2,0);
  \node[omDef] at (0.1,1.35) {$M$};
  \node[omThm, below] at (1.55,-0.05) {$\partial M$};
  \draw[->, thick, omExo] (1.415,1.415) -- (2.05,2.05);
  \node[omExo, right] at (2.02,2.12) {$\nu$};
  \draw[->, thick] (1.29,1.53) arc (50:75:2);
  \draw[->, thick, omExo] (-0.6,0) -- (-0.6,-0.65);
  \node[omExo, below] at (-0.6,-0.65) {$\nu$};
  \draw[->, thick] (-0.35,0) -- (0.55,0);
  \node[below] at (0.1,-0.14) {geïnduceerd};
\end{tikzpicture}

{\small De regel ``buitenwaartse normaal eerst'': zet in elk
randpunt de naar buiten wijzende vector $\nu$ vooraan; de bases
die haar tot een positief frame van $M$ aanvullen, oriënteren
$\partial M$. Voor een vlak gebied met de standaardoriëntatie is
dat de regel tegen de wijzers van de klok in van
Green--Riemann.}
\end{omfigure}

\begin{lemma}[Stokes op het halfvlak]
\label{lem:b3:forms:halfspace}
Zij $\eta$ een gladde $(k-1)$-vorm op $\R^k$ met \omterm{def:b3:topology:compact}{compacte} drager.
Dan is
\[
\int_{H^k}\dd\eta = \int_{\partial H^k}\eta,
\]
waarbij $H^k$ de standaardoriëntatie van $\R^k$ draagt en
$\partial H^k = \{u_k = 0\} \cong \R^{k-1}$ de geïnduceerde, die
$(-1)^k$ maal de standaardoriëntatie van $\R^{k-1}$ is.
\end{lemma}

\begin{proof}
Eerst de boekhouding van de \omterm{def:b3:forms:orientation}{oriëntatie}: de buitenwaartse normaal
langs $\partial H^k$ is $-e_k$, en
\[
\det(-e_k, e_1, \dots, e_{k-1}) =
-\det(e_k, e_1, \dots, e_{k-1}) = -(-1)^{k-1}\det(e_1,
\dots, e_k) = (-1)^k :
\]
het frame $(e_1, \dots, e_{k-1})$ van $\partial H^k$ is precies
positief voor de geïnduceerde \omterm{def:b3:forms:orientation}{oriëntatie} wanneer $k$ even is,
vandaar de vermelde vergelijking. Neem wegens de lineariteit
$\eta = f\,\dd u_1\wedge\dots\wedge\widehat{\dd
u_i}\wedge\dots\wedge\dd u_k$ met $f \in \mathcal
C^\infty_c(\R^k)$; dan is $\dd\eta =
(-1)^{i-1}\,\partial_if\,\dd u_1\wedge\dots\wedge\dd u_k$ ($\dd
u_i$ naar haar plaats verhuizen kost $i - 1$ verwisselingen).
Twee gevallen, beide met Tonelli--Fubini
(\cref{thm:b3:product:fubini}) en de hoofdstelling van de
integraalrekening in één veranderlijke.

\emph{Geval $i < k$.} Eerst in $u_i$ over $\R$ integreren:
$\int_\R\partial_if\,\dd u_i = 0$ (\omterm{def:b3:topology:compact}{compacte} drager), dus
$\int_{H^k}\dd\eta = 0$. En de beperking van $\eta$ tot $\{u_k =
0\}$ bevat de factor $\dd u_k$, die tot $0$ beperkt ($u_k$ is
daar constant): dus ook $\int_{\partial H^k}\eta = 0$.

\emph{Geval $i = k$.} Eerst in $u_k$ over $\intco0\infty$
integreren:
\begin{align*}
\int_{H^k}\dd\eta &= (-1)^{k-1}\int_{\R^{k-1}}
\Bigl(\int_0^\infty\partial_kf\,\dd u_k\Bigr)\dd u' \\
&= (-1)^{k-1}\int_{\R^{k-1}}\bigl(0 - f(u', 0)\bigr)\dd u'
= (-1)^k\int_{\R^{k-1}}f(u',0)\,\dd u' .
\end{align*}
Aan de randkant beperkt $\eta$ tot $f(u', 0)\,\dd
u_1\wedge\dots\wedge\dd u_{k-1}$, en omdat de geïnduceerde
\omterm{def:b3:forms:orientation}{oriëntatie} $(-1)^k$ maal de standaardoriëntatie is, is
$\int_{\partial H^k}\eta = (-1)^k\int_{\R^{k-1}}f(u',0)\,\dd u'$.
De twee leden vallen samen.
\end{proof}

\begin{theorem}[Stokes]\label{thm:b3:forms:stokes}
Zij $M \subseteq \R^n$ een \omterm{def:b3:topology:compact}{compacte} georiënteerde
$k$-deelvariëteit met \omterm{def:b3:forms:boundary}{rand}, waarbij $\partial M$ de geïnduceerde
\omterm{def:b3:forms:orientation}{oriëntatie} draagt, en zij $\omega$ een gladde $(k-1)$-vorm op een
\omterm{def:b3:topology:topology}{omgeving} van $M$. Dan is\index{stelling van Stokes}
\[
\int_M\dd\omega = \int_{\partial M}\omega .
\]
In het bijzonder is $\int_M\dd\omega = 0$ als $\partial M =
\varnothing$.
\end{theorem}

\begin{proof}
Overdek de \omterm{def:b3:topology:compact}{compacte} $M$ met eindig veel beelden van directe
kaarten (inwendig of \omterm{def:b3:forms:boundary}{rand}), en neem een \omterm{lem:b3:forms:partition}{partitie van de eenheid}
$(\chi_i)$ ondergeschikt aan de bijbehorende \omterm{def:b3:topology:topology}{open verzamelingen}
$W_i$ van $\R^n$ (\cref{lem:b3:forms:partition}), met
$\sum\chi_i = 1$ op een \omterm{def:b3:topology:topology}{omgeving} van $M$. Op die \omterm{def:b3:topology:topology}{omgeving} is
$\dd(\sum\chi_i) = 0$, dus
\[
\int_M\dd\omega = \sum_i\int_M\dd(\chi_i\omega),
\qquad
\int_{\partial M}\omega =
\sum_i\int_{\partial M}\chi_i\omega :
\]
beide leden zijn additief, en het volstaat de stelling te
bewijzen voor een vorm met drager in het beeld van één kaart.

\emph{Inwendige kaart.} Is $\operatorname{supp}\omega \cap M
\subseteq \gamma(V)$ met $V$ open in $\R^k$: breid $\eta =
\gamma^*\omega$ met nul uit tot $\R^k$ (glad, met \omterm{def:b3:topology:compact}{compacte} drager
in $V$) en pas \cref{lem:b3:forms:halfspace} toe met de drager
weg van $\partial H^k$ (verschuif $V$ het open bovenste halfvlak
in --- of herhaal eenvoudig de berekening van het geval $i < k$
over heel $\R^k$): dan is $\int_M\dd\omega = \int_{\R^k}\dd\eta =
0$, en $\int_{\partial M}\omega = 0$ omdat $\omega$ nabij
$\partial M$ verdwijnt.

\emph{Randkaart.} Is $\operatorname{supp}\omega\cap M \subseteq
\gamma(V \cap H^k)$: met $\eta = \gamma^*\omega$ met nul
uitgebreid is $\gamma^*(\dd\omega) = \dd\eta$
(\cref{thm:b3:forms:pullback}(b)), dus volgens
\cref{def:b3:forms:integral}:
\[
\int_M\dd\omega = \int_{H^k}\dd\eta
\overset{\text{\cref{lem:b3:forms:halfspace}}}{=}
\int_{\partial H^k}\eta .
\]
Rest nog het rechterlid met $\int_{\partial M}\omega$ te
identificeren. De \omterm{def:b3:forms:boundary}{rand} wordt geparametriseerd door $\beta(u') =
\gamma(u', 0)$, en $\beta^*\omega$ is de beperking van $\eta$ tot
$\{u_k = 0\}$ (de \omterm{def:b3:forms:pullback}{terugtrekking} onder de inclusie $u' \mapsto
(u', 0)$, samengesteld met $\gamma$). De vergelijking van de
\omterm{def:b3:forms:orientation}{oriëntaties} geeft aan beide kanten dezelfde $(-1)^k$: het frame
$(\partial_1\beta, \dots, \partial_{k-1}\beta)$ zit in de
geïnduceerde \omterm{def:b3:forms:orientation}{oriëntatie} van $\partial M$ met het teken
$\det(-e_k, e_1, \dots, e_{k-1}) = (-1)^k$, in de kaart berekend
(de buitenwaartse vector trekt terug naar $-e_k$), en dat is
precies het teken dat de geïnduceerde \omterm{def:b3:forms:orientation}{oriëntatie} van $\partial
H^k$ met de standaard van $\R^{k-1}$ verbindt
(\cref{lem:b3:forms:halfspace}). De twee tekenafspraken heffen
elkaar op: $\int_{\partial H^k}\eta = \int_{\partial M}\omega$.
\end{proof}

\begin{example}[De klassieke stellingen]
\label{ex:b3:forms:classical}
Zij $D \subseteq \R^2$ een \omterm{def:b3:topology:compact}{compact} gebied met gladde \omterm{def:b3:forms:boundary}{rand}, met de
standaardoriëntatie. Voor $\omega = P\,\dd x + Q\,\dd y$ is
$\dd\omega = \bigl(\partial_xQ - \partial_yP\bigr)\dd x\wedge\dd
y$, en luidt Stokes
\[
\int_D\Bigl(\frac{\partial Q}{\partial x} -
\frac{\partial P}{\partial y}\Bigr)\dd x\,\dd y
= \oint_{\partial D}P\,\dd x + Q\,\dd y :
\]
\emph{Green--Riemann}, bewezen voor elementaire gebieden in het
volume van bachelorjaar 2 en nu in natuurlijke algemeenheid. In
$\R^3$ geeft Stokes toegepast op de fluxvorm van een vectorveld
op een \omterm{def:b3:topology:compact}{compact} gebied de \emph{divergentiestelling}
$\int_\Omega\operatorname{div}F = \int_{\partial\Omega}\langle F,
\nu\rangle\,\dd S$, en toegepast op een $1$-vorm op een oppervlak
met \omterm{def:b3:forms:boundary}{rand} de klassieke rotatiestelling van
\emph{Kelvin--Stokes}; \cref{exo:b3:forms:7} schrijft beide
woordenboeken uit.
\end{example}

\section{Het windingsgetal}\label{sec:b3:forms:winding}

\begin{definition}\label{def:b3:forms:winding}
Zij $\gamma\colon\intcc01\to\R^2\setminus\{a\}$ een gladde
\omterm{def:b3:forms:closedexact}{gesloten} kromme. Haar \emph{windingsgetal}\index{windingsgetal}
rond $a$ is
\[
\operatorname{Ind}_\gamma(a) =
\frac1{2\pi}\int_\gamma\omega_\theta^a,
\qquad
\omega_\theta^a = \frac{(x - a_1)\,\dd y - (y - a_2)\,\dd
x}{(x - a_1)^2 + (y - a_2)^2} .
\]
\end{definition}

\begin{proposition}\label{prop:b3:forms:windinginteger}
$\operatorname{Ind}_\gamma(a) \in \Z$; als functie van $a$ is ze
constant op elke \omterm{def:b3:topology:components}{samenhangscomponent} van
$\R^2\setminus\gamma(\intcc01)$ en nul op de onbegrensde
component. Voor $\gamma(t) = a + r(\cos2\pi t, \sin2\pi t)$ is
$\operatorname{Ind}_\gamma(a) = 1$.
\end{proposition}

\begin{proof}
Neem $a = 0$ en schrijf $\gamma = (x, y)$, $\rho = \norm\gamma >
0$ en $\vartheta(t) = \int_0^t\gamma^*\omega_\theta$, zodat
$\vartheta' = \frac{xy' - yx'}{x^2 + y^2}$. Zij in complexe
notatie $u(t) = \gamma(t)\,\rho(t)^{-1}\eu^{-\iu\vartheta(t)}$;
dan is $\abs u = 1$ en geeft een rechtstreekse berekening
\[
\frac{u'}{u} = \frac{\gamma'}{\gamma} - \frac{\rho'}{\rho}
- \iu\vartheta'
= \frac{\bar\gamma\gamma'}{\abs\gamma^2} -
\frac{\rho'}{\rho} - \iu\,\frac{xy' -
yx'}{\rho^2}
= \frac{(xx' + yy')}{\rho^2} - \frac{\rho'}{\rho} = 0,
\]
want $\bar\gamma\gamma' = (xx' + yy') + \iu(xy' - yx')$ en
$\rho\rho' = xx' + yy'$. Dus is $u$ constant: $\gamma(t) =
\rho(t)\,c\,\eu^{\iu\vartheta(t)}$ met $\abs c = 1$, en
$\gamma(1) = \gamma(0)$ met $\rho(1) = \rho(0)$ dwingt
$\eu^{\iu\vartheta(1)} = \eu^{\iu\vartheta(0)} = 1$ af:
$\vartheta(1) \in 2\pi\Z$, dat wil zeggen
$\operatorname{Ind}_\gamma(0) \in \Z$. Als functie van $a$ op het
open complement van de \omterm{def:b3:topology:compact}{compacte} kromme is de definiërende
integraal \omterm{def:b3:topology:continuity}{continu} (\cref{thm:b3:lebesgue:paramcont}, met
dominantie op een \omterm{def:b3:topology:topology}{omgeving} van elke $a$); een \omterm{def:b3:topology:continuity}{continue} functie
met gehele waarden is lokaal constant, dus constant op de
componenten. Voor grote $\norm a$ is de integrand uniform in $t$
$O(1/\norm a)$, dus gaat de index naar $0$ en verdwijnt ze op de
onbegrensde component. Voor de cirkel:
$\gamma^*\omega_\theta^a = 2\pi\,\dd t$ rechtstreeks.
\end{proof}

\begin{remark}
Onder $\C \cong \R^2$ is $\frac{\dd z}{z} = \frac{x\,\dd x +
y\,\dd y}{x^2 + y^2} + \iu\,\omega_\theta$, dus
$\operatorname{Ind}_\gamma(a) =
\frac1{2\iu\pi}\oint_\gamma\frac{\dd z}{z - a}$: dat is de index
van \cref{ch:b3:residues}, en
\cref{prop:b3:forms:windinginteger} herbewijst haar
geheeltalligheid en lokale constantheid met reële middelen ---
de topologische helft van de residustelling, nu staand op
Stokes.
\end{remark}

\begin{proposition}\label{prop:b3:forms:exactloop}
Is $\omega = \dd f$ \omterm{def:b3:forms:closedexact}{exact} op de open $U$ en is
$\gamma\colon\intcc01\to U$ een \omterm{def:b3:forms:closedexact}{gesloten} kromme, dan is
$\int_\gamma\omega = 0$.
\end{proposition}

\begin{proof}
$\int_\gamma\dd f = \int_0^1(f\circ\gamma)'(t)\,\dd t =
f(\gamma(1)) - f(\gamma(0)) = 0$ --- de kettingregel
identificeert $\gamma^*(\dd f)$ met $(f\circ\gamma)'\,\dd t$.
\end{proof}

\begin{method}\label{met:b3:forms:method}
Rekenen met vormen: (1) mechaniseer --- \omterm{def:b3:forms:wedge}{uitwendige producten}
herordenen met tekens, $\dd$ differentieert de coëfficiënten naar
verse $\dd x_j$'s, \omterm{def:b3:forms:pullback}{terugtrekkingen} substitueren; vertrouw de
algebra, ze codeert elke jacobiaan. (2) Om een vorm over een
deelvariëteit te integreren: parametriseer direct, trek terug,
integreer de coëfficiënt; de \omterm{def:b3:forms:orientation}{oriëntatie} is de enige val ---
controleer één frame. (3) Om een integraalidentiteit te bewijzen,
zoek naar een vorm van Stokes: is de integrand \omterm{def:b3:forms:closedexact}{exact}? is het
gebied een \omterm{def:b3:forms:boundary}{rand}? (4) Om integralen over twee ``evenwijdige''
deelvariëteiten te vergelijken, pas Stokes toe op het gebied
ertussen (het vervormingsargument, \cref{exo:b3:forms:10}). (5)
Een integraal $\neq 0$ van een gesloten vorm certificeert een
topologische obstructie --- geen primitieve, geen retractie, geen
nulpuntvrije uitbreiding: zo doodt de weekendopgave de retracties
van de bal.
\end{method}

\section{Oefeningen}

\begin{exercise}[$\star$]\label{exo:b3:forms:1}
Zij op $\R^3$ $\omega = x\,\dd y - z\,\dd x$ en $\eta = \dd x +
y\,\dd z$. Bereken $\omega\wedge\eta$, $\dd\omega$, $\dd\eta$ en
$\dd(\omega\wedge\eta)$, en ga de gegradeerde regel van Leibniz
op dit voorbeeld na.
\end{exercise}

\begin{exercise}[$\star$]\label{exo:b3:forms:2}
Identificeer op $\R^3$ de drie gedaanten van $\dd$: voor $f \in
\Omega^0$ is $\dd f \leftrightarrow \nabla f$; voor de arbeidsvorm
$\omega_F = F_1\dd x + F_2\dd y + F_3\dd z$ is $\dd\omega_F
\leftrightarrow \operatorname{curl}F$; voor de fluxvorm $\sigma_F
= F_1\,\dd y\wedge\dd z + F_2\,\dd z\wedge\dd x + F_3\,\dd
x\wedge\dd y$ is $\dd\sigma_F \leftrightarrow
\operatorname{div}F$. Leid uit $\dd^2 = 0$ de identiteiten
$\operatorname{curl}\nabla f = 0$ en
$\operatorname{div}\operatorname{curl}F = 0$ af.
\end{exercise}

\begin{exercise}[$\star\star$]\label{exo:b3:forms:3}
Beslis voor elke $1$-vorm of ze \omterm{def:b3:forms:closedexact}{gesloten} is, of ze \omterm{def:b3:forms:closedexact}{exact} is op
haar domein, en bereken een primitieve wanneer die bestaat:
(a) $(2xy + z^2)\,\dd x + x^2\,\dd y + 2xz\,\dd z$ op $\R^3$;
(b) $\dfrac{x\,\dd x + y\,\dd y}{x^2 + y^2}$ op
$\R^2\setminus\{0\}$;
(c) $\dfrac{-y\,\dd x + x\,\dd y}{x^2 + y^2}$ op het halfvlak
$\{x > 0\}$.
\end{exercise}

\begin{exercise}[$\star\star$]\label{exo:b3:forms:4}
(De hoekvorm) Ga na dat $\omega_\theta$
(\cref{ex:b3:forms:angular}) \omterm{def:b3:forms:closedexact}{gesloten} is; bereken
$\int_\gamma\omega_\theta$ voor $\gamma$ de cirkel met straal $r$
rond $0$; besluit dat $\omega_\theta$ niet \omterm{def:b3:forms:closedexact}{exact} is op
$\R^2\setminus\{0\}$, en dat $\R^2\setminus\{0\}$ ten opzichte
van geen van haar punten stervormig is (twee wegen: via
\cref{thm:b3:forms:poincare}, en rechtstreeks uit de
meetkunde).
\end{exercise}

\begin{exercise}[$\star\star$]\label{exo:b3:forms:5}
(Oppervlaktevorm van een hyperoppervlak) Zij $M \subseteq \R^n$
een \omterm{def:b3:topology:compact}{compact} georiënteerd hyperoppervlak waarvan de \omterm{def:b3:forms:orientation}{oriëntatie}
gegeven wordt door een normaalveld $\nu$ van norm $1$ ($(w_1,
\dots, w_{n-1})$ is positief dan en slechts dan als $(\nu, w_1,
\dots, w_{n-1})$ positief is in $\R^n$). Toon aan dat de
$(n-1)$-vorm $\sigma_\nu =
\sum_i(-1)^{i-1}\nu_i\,\dd x_1\wedge\dots\wedge\widehat{\dd
x_i}\wedge\dots\wedge\dd x_n$ op $M$ beperkt tot de
\emph{oppervlaktevorm}: voor een directe parametrisering $\gamma$
is $\gamma^*\sigma_\nu = \sqrt{\det G}\,\dd u_1\wedge\dots
\wedge\dd u_{n-1}$, waarbij $G = ({}^t D\gamma)(D\gamma)$ de
grammatrix is. \emph{(Merk op dat $(\gamma^*\sigma_\nu)(e_1,
\dots, e_{n-1}) = \det(\nu, \partial_1\gamma, \dots,
\partial_{n-1}\gamma)$ en kwadrateer die determinant.)} Bereken
$\int_{S^2}\sigma_\nu = 4\pi$.
\end{exercise}

\begin{exercise}[$\star\star$]\label{exo:b3:forms:6}
Bereken $\int_{S^2}\omega$ voor $\omega = x\,\dd y\wedge\dd z +
y\,\dd z\wedge\dd x + z\,\dd x\wedge\dd y$: (a) rechtstreeks in
bolcoördinaten; (b) met Stokes op de eenheidsbal. Leid
$\operatorname{vol}(B^3) = \frac{4\pi}3$ af uit de oppervlakte
van $S^2$, en veralgemeen: $n\operatorname{vol}(B^n) =
\operatorname{area}(S^{n-1})$, in overeenstemming met
\cref{thm:b3:product:ballvolume}.
\end{exercise}

\begin{exercise}[$\star\star$]\label{exo:b3:forms:7}
(De woordenboeken) Leid zorgvuldig uit
\cref{thm:b3:forms:stokes} af: (a) de divergentiestelling in
$\R^3$ (combineer \cref{exo:b3:forms:2} en
\cref{exo:b3:forms:5}); (b) de stelling van Kelvin--Stokes
$\int_S\langle\operatorname{curl}F, \nu\rangle\,\dd S =
\oint_{\partial S}\langle F, \tau\rangle\,\dd\ell$ voor een
\omterm{def:b3:topology:compact}{compact} georiënteerd oppervlak met \omterm{def:b3:forms:boundary}{rand} in $\R^3$. Controleer dat
de oriëntatieafspraken overeenstemmen op de bovenste halve bol
begrensd door de evenaar.
\end{exercise}

\begin{exercise}[$\star\star$]\label{exo:b3:forms:8}
(Identiteiten van Green) Bewijs voor $u, v$ glad op een \omterm{def:b3:topology:topology}{omgeving}
van een \omterm{def:b3:topology:compact}{compact} gebied $\Omega \subseteq \R^n$ met gladde \omterm{def:b3:forms:boundary}{rand}
dat
\[
\int_\Omega\bigl(u\,\Delta v + \langle\nabla u, \nabla
v\rangle\bigr) = \int_{\partial\Omega}u\,\partial_\nu
v\,\dd S,
\qquad
\int_\Omega\bigl(u\,\Delta v - v\,\Delta u\bigr) =
\int_{\partial\Omega}\bigl(u\,\partial_\nu v -
v\,\partial_\nu u\bigr)\dd S .
\]
Leid af: een \omterm{prop:b3:conformal:harmonicholo}{harmonische functie} op $\Omega$ die op
$\partial\Omega$ verdwijnt, verdwijnt identiek, en twee
\omterm{prop:b3:conformal:harmonicholo}{harmonische functies} met dezelfde randwaarden vallen samen --- de
uniciteit in het dirichletprobleem van
\cref{ch:b3:conformal}.
\end{exercise}

\begin{exercise}[$\star\star$]\label{exo:b3:forms:9}
(Variabelensubstitutie, georiënteerde vorm) Zij $\varphi\colon U
\to V$ een diffeomorfisme van open delen van $\R^n$ met $\det
D\varphi > 0$, en $f$ \omterm{def:b3:topology:continuity}{continu} met \omterm{def:b3:topology:compact}{compacte} drager in $V$. Toon
aan dat de terugtrekkingsidentiteit $\int_U\varphi^*(f\,\dd
x_1\wedge\dots\wedge\dd x_n) = \int_Vf\,\dd x_1\wedge\dots\wedge
\dd x_n$ gelijkwaardig is met de stelling over de
variabelensubstitutie (\cref{thm:b3:product:changeofvar}) voor
zulke $\varphi$, en leg precies uit waar de absolute waarde op de
jacobiaan is gebleven.
\end{exercise}

\begin{exercise}[$\star\star\star$]\label{exo:b3:forms:10}
(Vervorming) Zij $\omega$ een \emph{\omterm{def:b3:forms:closedexact}{gesloten}} $2$-vorm op
$\R^3\setminus\{0\}$ en $S_r$ de sfeer met straal $r$ rond $0$.
Toon aan dat $\int_{S_r}\omega$ niet van $r > 0$ afhangt
\emph{(pas Stokes toe op de schil tussen twee stralen; let op de
twee randoriëntaties)}. Pas dat toe op de \emph{ruimtehoekvorm}
\[
\omega = \frac{x\,\dd y\wedge\dd z + y\,\dd z\wedge\dd x +
z\,\dd x\wedge\dd y}{(x^2 + y^2 + z^2)^{3/2}} :
\]
ga na dat ze \omterm{def:b3:forms:closedexact}{gesloten} is, bereken $\int_{S_r}\omega = 4\pi$, en
besluit dat ze \omterm{def:b3:forms:closedexact}{gesloten} maar niet \omterm{def:b3:forms:closedexact}{exact} is op
$\R^3\setminus\{0\}$ --- het tweedimensionale zusje van
$\omega_\theta$, en de meetkundige inhoud van de wet van Gauss in
de elektrostatica.
\end{exercise}

\begin{exercise}[$\star\star$]\label{exo:b3:forms:11}
Zij $\gamma$ een gladde \omterm{def:b3:forms:closedexact}{gesloten} kromme in $\R^2\setminus\{0\}$
met $n = \operatorname{Ind}_\gamma(0)$. Toon aan dat
$\int_\gamma\omega = n\int_{S^1}\omega$ voor \emph{elke} \omterm{def:b3:forms:closedexact}{gesloten}
$1$-vorm $\omega$ op $\R^2\setminus\{0\}$ \emph{(schrijf $\omega
= c\,\omega_\theta + \dd f$ volgens \cref{exo:b3:forms:12})}.
Interpretatie: op het geperforeerde vlak is het \omterm{def:b3:forms:winding}{windingsgetal} de
enige obstructie voor het verdwijnen van de perioden.
\end{exercise}

\begin{exercise}[$\star\star\star$]\label{exo:b3:forms:12}
(Eerste berekening van de Rham) Toon aan dat elke \omterm{def:b3:forms:closedexact}{gesloten}
$1$-vorm $\omega$ op $U = \R^2\setminus\{0\}$ op unieke wijze
\[
\omega = c\,\omega_\theta + \dd f,
\qquad c = \frac1{2\pi}\int_{S^1}\omega,\quad f \in
\mathcal C^\infty(U) :
\]
definieer $f(p)$ door $\omega - c\,\omega_\theta$ te integreren
langs een weg van $(1,0)$ naar $p$ (eerst een radiaal stuk, dan
een cirkelboog), toon aan dat het resultaat niet van de keuzes
afhangt juist omdat de periode over $S^1$ verdwijnt, en ga $\dd f
= \omega - c\,\omega_\theta$ na. Besluit:
$H^1(\R^2\setminus\{0\}) \cong \R$, voortgebracht door de
hoekvorm.
\end{exercise}

\section{Probleem: de vastepuntstelling van Brouwer}

\begin{problem}[{Weekendopgave --- geen retractie, geen
ontsnapping}]\label{pb:b3:forms:1}
De stelling van Brouwer zegt dat \emph{elke \omterm{def:b3:topology:continuity}{continue afbeelding}
van de \omterm{def:b3:forms:closedexact}{gesloten} eenheidsbal $\bar B = \bar B^n \subseteq \R^n$ in
zichzelf een vast punt heeft} --- een van de grote stellingen van
de wiskunde, met gevolgen van de speltheorie (nash-evenwichten)
tot de matrixanalyse. Het bewijs met \omterm{def:b3:forms:diffform}{differentiaalvormen} is het
netste dat bekend is: Stokes toont dat de sfeer geen retract van
de bal is, en de rest volgt. Overal is $S = S^{n-1} = \partial\bar
B$ met $n \geq 2$, en
\[
\sigma = \sum_{i=1}^n(-1)^{i-1}x_i\,
\dd x_1\wedge\dots\wedge\widehat{\dd x_i}\wedge\dots\wedge
\dd x_n
\]
(het dakje schrapt de factor).

\medskip
\noindent\textbf{Deel I --- Het meetinstrument.}
\begin{enumerate}
  \item Bereken $\dd\sigma$, en leid met Stokes
        (\cref{thm:b3:forms:stokes}) af dat $\int_S\sigma =
        n\operatorname{vol}(\bar B) > 0$, waarbij $S$ de
        randoriëntatie van de bal draagt.
  \item Identificeer voor $n = 2$ en $n = 3$ de beperking van
        $\sigma$ tot $S$ met de booglengte- en de
        oppervlaktevorm (\cref{exo:b3:forms:5} met $\nu(x) = x$)
        en herbereken $\int_S\sigma$ rechtstreeks.
  \item Zij $W \subseteq \R^N$ open en $\varphi\colon W \to \R^n$
        glad met $\norm{\varphi(x)} = 1$ voor alle $x \in W$.
        Toon aan dat $\varphi^*(\dd\sigma) = 0$.
        \emph{(Differentieer $\norm\varphi^2 = 1$: het beeld van
        $D\varphi(x)$ ligt in het hypervlak $\varphi(x)^\perp$,
        van dimensie $n - 1$; pas dan
        \cref{prop:b3:forms:linearpullback}(2) toe.)}
  \item Waar zit de fout in het volgende ``bewijs'' dat
        $\int_S\sigma = 0$: ``$S$ is \omterm{def:b3:topology:compact}{compact} zonder \omterm{def:b3:forms:boundary}{rand}, en
        $\sigma$ beperkt tot $S$ is er een vorm van topgraad, dus
        \omterm{def:b3:forms:closedexact}{gesloten}, dus $\int_S\sigma = \int_S\dd(\text{iets}) = 0$
        volgens Stokes''? (Wijs het foute woord aan.)
  \item Leg in één alinea de strategie van Deel II uit: wat zal
        er worden geïntegreerd, waarover, en waar de tegenspraak
        vandaan zal komen.
\end{enumerate}

\medskip
\noindent\textbf{Deel II --- Geen gladde retractie.} Stel, om tot
een tegenspraak te komen, dat $r$ een \emph{gladde retractie} van
de bal op haar sfeer is: $r$ is glad op een \omterm{def:b3:topology:topology}{omgeving} van $\bar
B$, $r(\bar B) \subseteq S$, en $r(x) = x$ voor alle $x \in S$.
\begin{enumerate}[resume]
  \item Verantwoord $\int_Sr^*\sigma = \int_S\sigma$ \emph{(op
        $S$ beperkt $r$ tot de identiteit: is $\gamma$ een
        directe parametrisering van een stuk van $S$, dan is
        $r\circ\gamma = \gamma$)}.
  \item Toon met Stokes op $\bar B$ aan dat $\int_Sr^*\sigma =
        \int_{\bar B}\dd(r^*\sigma)$.
  \item Toon aan dat $\dd(r^*\sigma) = r^*(\dd\sigma) = 0$ (vraag
        3 toegepast op $\varphi = r$), en besluit: \emph{er
        bestaat geen gladde retractie $\bar B \to S$.}
  \item Handel het uitgesloten geval $n = 1$ met de hand af: toon
        rechtstreeks aan dat geen enkele \omterm{def:b3:topology:continuity}{continue afbeelding}
        $\intcc{-1}1 \to \{-1, 1\}$ beide eindpunten vastlaat, en
        noem de stelling die je gebruikte.
\end{enumerate}

\medskip
\noindent\textbf{Deel III --- Gladde Brouwer.} Zij $g$ glad op
een \omterm{def:b3:topology:topology}{omgeving} van $\bar B$ met $g(\bar B) \subseteq \bar B$ en
\emph{geen} vast punt in $\bar B$.
\begin{enumerate}[resume]
  \item Toon aan dat $\delta = \min_{x\in\bar B}\norm{g(x) - x} >
        0$.
  \item Zij voor $x \in \bar B$ $u(x) = \frac{x -
        g(x)}{\norm{x - g(x)}}$ en zij $r(x) = x + t(x)\,u(x)$
        het snijpunt van de straal $\{x + tu(x) : t \geq 0\}$ met
        $S$. Los de tweedegraadsvergelijking $\norm{x + tu}^2 =
        1$ op en verkrijg
        \[
        t(x) = -\langle x, u(x)\rangle +
        \sqrt{1 - \norm x^2 + \langle x, u(x)\rangle^2}
        \;\geq\; 0 .
        \]
  \item Toon aan dat de uitdrukking onder de wortel op $\bar B$
        \emph{strikt} positief is: is $1 - \norm x^2 + \langle x,
        u(x)\rangle^2 = 0$, dan is $\norm x = 1$ en $\langle x,
        u(x)\rangle = 0$, dat wil zeggen $\langle x, x -
        g(x)\rangle = 0$, oftewel $\langle x, g(x)\rangle = 1$;
        met Cauchy--Schwarz en $\norm x = 1$, $\norm{g(x)} \leq
        1$ dwingt dat $g(x) = x$ af --- uitgesloten. Leid af dat
        $r$ glad is op een \omterm{def:b3:topology:topology}{omgeving} van $\bar B$.
  \item Toon aan dat $r(\bar B) \subseteq S$ en $r(x) = x$ voor
        $x \in S$ \emph{(controleer voor $\norm x = 1$ dat $t(x)
        = 0$ met $\langle x, u(x)\rangle \geq 0$, wat zelf volgt
        uit $\langle x, x - g(x)\rangle = 1 - \langle x,
        g(x)\rangle \geq 0$)}. Besluit met Deel II: \emph{elke
        gladde zelfafbeelding van $\bar B$ heeft een vast
        punt.}
\end{enumerate}

\medskip
\noindent\textbf{Deel IV --- \omterm{def:b3:topology:continuity}{Continue} Brouwer.} Zij
$f\colon\bar B\to\bar B$ \omterm{def:b3:topology:continuity}{continu} zonder vast punt.
\begin{enumerate}[resume]
  \item Toon aan dat $\varepsilon = \min_{\bar B}\norm{f - x} >
        0$, en geef een \emph{veeltermafbeelding} $p\colon\R^n\to
        \R^n$ met $\sup_{\bar B}\norm{p - f} < \varepsilon/2$
        (Stone--Weierstrass,
        \cref{thm:b3:complete:stoneweierstrass}, coördinaat per
        coördinaat --- verantwoord de overgang van de scalaire
        naar de vectoriële benadering).
  \item De afbeelding $p$ kan de bal verlaten; stel $g =
        \frac{p}{1 + \varepsilon/2}$. Toon aan dat $g(\bar B)
        \subseteq \bar B$ en $\sup_{\bar B}\norm{g - f} <
        \varepsilon$.
  \item Leid een tegenspraak af met Deel III en besluit:
        \emph{elke \omterm{def:b3:topology:continuity}{continue afbeelding} $\bar B^n \to \bar B^n$
        heeft een vast punt.}
  \item Toon met voorbeelden aan dat de stelling faalt op: de
        open bal; de sfeer $S$; een \omterm{def:b3:forms:closedexact}{gesloten} ring. Welke
        eigenschap van $\bar B$ verliest elk tegenvoorbeeld?
\end{enumerate}

\medskip
\noindent\textbf{Deel V --- Dividenden.}
\begin{enumerate}[resume]
  \item (Perron--Frobenius, bestaan) Zij $A$ een $n\times
        n$-matrix met alle elementen $> 0$, en $\Delta = \{x \in
        \R^n : x_i \geq 0,\ \sum x_i = 1\}$. Toon aan dat de
        afbeelding $x \mapsto Ax/\norm{Ax}_1$ goed gedefinieerd
        en \omterm{def:b3:topology:continuity}{continu} is op $\Delta$, dat $\Delta$ \omterm{def:b3:topology:continuity}{homeomorf} is met
        een \omterm{def:b3:forms:closedexact}{gesloten} bal van $\R^{n-1}$ (het radiale
        \omterm{def:b3:topology:continuity}{homeomorfisme} vanuit een convexe \omterm{def:b3:topology:compact}{compacte} verzameling met
        niet-leeg inwendige in haar affiene omhulsel), en besluit
        dat $A$ een eigenvector met strikt positieve componenten
        en een eigenwaarde $> 0$ heeft.
  \item Leid af dat elke stochastische matrix met positieve
        elementen (kolommen met som $1$) een stationaire
        kansvector $\pi = A\pi$ heeft --- de vector van het
        pageranktype. (De uniciteit geldt eveneens, maar vergt
        ander gereedschap.)
  \item (Harige bal, opzet) Zij $v$ glad op een \omterm{def:b3:topology:topology}{omgeving} van $S =
        S^{n-1}$ met $\langle v(x), x\rangle = 0$ en
        $\norm{v(x)} = 1$ voor $x \in S$ (een \emph{raakveld van
        norm $1$}). Stel voor $t \in \R$ $F_t(x) = x + t\,v(x)$.
        Toon aan dat $\norm{F_t(x)} = \sqrt{1 + t^2}$ op $S$:
        $F_t$ beeldt $S$ af in de sfeer $\sqrt{1+t^2}\,S$.
  \item Toon aan dat $P(t) = \int_SF_t^*\sigma$ een
        \emph{veelterm} in $t$ is \emph{(elke coëfficiëntfunctie
        van $F_t^*\sigma$ in een kaart is polynomiaal in $t$, met
        coëfficiënten die glad zijn in de kaartveranderlijke; de
        integratie is lineair)}.
  \item Toon aan dat $F_t$ voor kleine $\abs t$ een
        diffeomorfisme is van $S$ op $\sqrt{1+t^2}\,S$: injectief
        voor $t\operatorname{Lip}(v) < 1$; een lokaal
        diffeomorfisme volgens de inverse-functiestelling
        (\cref{thm:b3:submanifolds:ift} in kaarten); en het beeld
        open en \omterm{def:b3:forms:closedexact}{gesloten} in de \omterm{def:b3:topology:connected}{samenhangende} doelsfeer. Leid met
        \cref{lem:b3:forms:consistency} en de schaling
        $\sigma_{\lambda x} = \lambda^{n}\,\sigma_x$ onder $x
        \mapsto \lambda x$ (ga die na) af dat
        \[
        P(t) = \pm(1 + t^2)^{n/2}\int_S\sigma,
        \qquad\text{met het teken } + \text{ voor kleine } t
        \]
        (de \omterm{def:b3:forms:orientation}{oriëntatie} blijft behouden wegens de \omterm{def:b3:topology:continuity}{continuïteit}
        vanaf $t = 0$).
  \item Besluit (Milnor): is $n$ \emph{oneven}, dan is $(1 +
        t^2)^{n/2}$ geen veelterm in $t$, en toch valt ze nabij
        $0$ samen met de veelterm $P(t)/\int_S\sigma$ ---
        tegenspraak. Bijgevolg dragen \emph{de
        evendimensionale sferen $S^{n-1}$ (met $n$ oneven) geen
        enkel raakveld van norm $1$}, en, na normeren en
        gladstrijken (convolueer componentsgewijs en projecteer
        --- verantwoord beide stappen), helemaal geen \omterm{def:b3:topology:continuity}{continu}
        raakveld zonder nulpunten: elke wind op aarde laat een
        windstil punt achter.

  \item (Oneven sferen laten zich kammen) Geef op $S^{2m-1}
        \subseteq \R^{2m} \cong \C^m$ een expliciet glad raakveld
        van norm $1$: $v(x) = \iu x$ in complexe notatie, dat wil
        zeggen
        \[
        v(x_1, y_1, \dots, x_m, y_m) =
        (-y_1, x_1, \dots, -y_m, x_m) .
        \]
        Ga de raking en de norm na, en besluit dat de
        pariteitsdichotomie van vraag 23 scherp is: een sfeer
        laat zich precies kammen wanneer haar dimensie oneven is.
        Waar breekt het veeltermargument van vraag 22 voor even
        $n$?
  \item (Nulpunten uit randgedrag) Zij $f \colon \bar B^n \to
        \R^n$ \omterm{def:b3:topology:continuity}{continu} met $\langle f(x), x\rangle \geq 0$ voor
        elke $x \in S^{n-1}$. Toon aan dat $f$ ergens in $\bar
        B^n$ verdwijnt. \emph{(Zo niet, dan beeldt $g(x) =
        -f(x)/\norm{f(x)}$ $\bar B$ \omterm{def:b3:topology:continuity}{continu} af in $S \subseteq
        \bar B$; pas Brouwer toe op $g$ en weerspreek de
        randhypothese.)} Leid het \emph{surjectiviteitscriterium}
        af: een \omterm{def:b3:topology:continuity}{continue} $F\colon\R^n\to\R^n$ met
        $\frac{\langle F(x), x\rangle}{\norm x} \to +\infty$ als
        $\norm x \to \infty$ is surjectief --- de
        eindigdimensionale voorvader van de
        coërciviteitsargumenten van de niet-lineaire analyse.
\end{enumerate}
\end{problem}
